\chapter{Project Context and Motivation}
\label{chap:project-context-motivation}

\section{Urban Isolation in Apartment Communities}
In modern metropolitan centers, high-density residential developments have paradoxically fostered social isolation. Neighbors live in physical proximity, separated only by thin concrete partitions, yet they often share no meaningful social capital or interpersonal trust. This social fragmentation is part of a broader decline in localized civic participation and community cohesion, a phenomenon well-documented by sociologists such as Robert Putnam in his seminal work on the erosion of social capital \cite{putnam2000bowling}. 

The consequences of this disconnect are not merely social, but highly practical. Without active communication channels, residents miss opportunities for mutual aid, shared resource allocation, building governance coordination, and safety collaboration. Simple daily challenges, such as finding a temporary parking spot, borrowing a household tool, or planning building-wide maintenance, become friction-filled tasks because there is no trusted hyper-local communication platform.

\section{Local Digital Networks and Mutual Aid}
To bridge this social gap, urban informatics researchers have investigated how digital communication networks can be designed to support physical communities. Keith Hampton and Barry Wellman's research on localized networks demonstrates that web-based platforms, when designed specifically for geographically bounded neighborhoods, act as catalysts for real-world social interaction rather than replacement channels \cite{hampton2003neighboring}.

Hyper-local systems such as Nextdoor \cite{nextdoor2024} or building-wide WhatsApp groups are commonly deployed, yet they exhibit significant limitations:
\begin{itemize}
    \item \textbf{Commercial Noise and Geographic Limitations:} Global platforms like Nextdoor are heavily commercialized and suffer from strict regional restrictions, making them unavailable in Romania and many other parts of the world. Additionally, they introduce commercial noise and pose privacy concerns that undermine local resident trust \cite{masden2014tensions}.
    \item \textbf{Informal Group Chat Chaos:} Instant messaging channels, such as WhatsApp or Telegram, lack structured categorization. Announcements, item lending requests, and event coordinates quickly become lost in continuous text streams, causing user fatigue.
    \item \textbf{Lack of Verification:} General social networks do not provide building-level verification. Anyone can join, making transactions like tool sharing or parking allocation risky.
\end{itemize}

NEST addresses these limitations by providing a structured, building-specific digital space. It is designed to turn online coordination into real-world cooperation.

\section{NEST Project Overview}
NEST is a hyper-local web application designed exclusively for residents of a single apartment building block. It provides a secure, structured space where residents can build trust, coordinate mutual aid, and organize community initiatives.

Rather than trying to replicate a broad social network, the system focuses on structured utility modules:
\begin{itemize}
    \item \textbf{Social Hub:} Houses the chronological Community Feed where verified residents share general announcements and building news.
    \item \textbf{Collaborative Consumption:} Operates the Shared Shed, enabling safe peer-to-peer household item lending and borrowing.
    \item \textbf{Spatial Optimization:} Integrates the Parking Sharing system, helping residents coordinate temporary parking slot access.
    \item \textbf{Physical Organizing:} Supports the Events module to facilitate real-world neighborhood meetings and gatherings.
\end{itemize}

My contribution to the NEST project is the frontend design and implementation, utilizing Angular 21, reactive Signals, and NgRx state management, alongside the integration of robust backend services, JWT verification middleware, database modeling with Prisma ORM, and the custom color CSS design token system.

\section{Project Objectives}
The NEST platform was developed with five core objectives:
\begin{enumerate}
    \item \textbf{Objective 1 (O1): Building Verification Security} -- Design and implement a secure, multi-stage registration system requiring administrative approval before granting residents access to building data.
    \item \textbf{Objective 2 (O2): Cohesive Design System} -- Create a unified CSS design token system that supports responsive layouts and instant light/dark theme transitions across all modules.
    \item \textbf{Objective 3 (O3): Collaborative Consumption (Shared Shed)} -- Implement a peer-to-peer tool-sharing catalog with conflict prevention and transaction history tracking.
    \item \textbf{Objective 4 (O4): Spatial Sharing (Parking)} -- Build a parking spot sharing system with auto-rejection logic for overlapping requests.
    \item \textbf{Objective 5 (O5): High-Performance Reactive Architecture} -- Structure state management using the NgRx Facade pattern and Angular Signals to ensure fast initial loads and fine-grained DOM rendering.
\end{enumerate}

\section{Original Contributions}
This thesis documents several technical contributions in localized web engineering:
\begin{itemize}
    \item \textbf{Adapted Facade-Signals Reactive Architecture:} Adapted the established NgRx Facade design pattern and integrated it with Angular's modern Signal-based reactivity system to build a state management architecture tailored to NEST's multi-module structure, isolating store mechanics from UI components while leveraging fine-grained template reactivity.
    \item \textbf{Dynamic CSS Theme Engine:} Constructed a design token system containing over 600 custom properties, achieving instant light-to-dark theme transitions without DOM re-renders.
    \item \textbf{Secure Multi-Stage Verification Flow:} Developed a routing and security guard pipeline that limits unverified users to onboarding views until approved by the building administrator.
    \item \textbf{Conflict-Resilient Rental \& Lending Logic:} Implemented backend transaction checks using Prisma ORM to prevent double-booking of parking slots and shared tools.
\end{itemize}

\section{Thesis Structure}
The remainder of this thesis is organized as follows. Chapter~2 reviews the theoretical background and state of the art, covering single-page application paradigms, reactive state management, CSS architecture, web security, and related community platforms. Chapter~3 presents the software analysis and design phase, detailing the requirements specification, system architecture, database modeling, and REST API contract definitions. Chapter~4 describes the implementation details, including the codebase organization, component structures, the NgRx Facade integration, core module logic, and security middleware. Chapter~5 discusses the testing and validation strategies, presenting build verification outputs, cross-browser compatibility checks, and usability considerations. Finally, Chapter~6 summarizes the project outcomes, identifies current engineering limitations, and proposes directions for future development.

\section{Declaration on the Use of AI Tools}
In compliance with the academic integrity guidelines of Babe\c{s}-Bolyai University (Faculty of Mathematics and Computer Science), I declare that generative Artificial Intelligence tools, specifically large language models, were utilized during the development and writing phases of this thesis. 

AI assistance was specifically applied for:
\begin{itemize}
    \item \textbf{Text Rephrasing:} Assisting in rewriting drafts of the thesis text to improve syntactic flow, clarity, and grammatical cohesion.
    \item \textbf{Technical Expression Enhancement:} Refining descriptions of complex software engineering concepts and architectural styles to align with professional, academic phrasing standards.
    \item \textbf{UI Design Generation:} The AI assistant integrated within Figma was used to help generate user interface design mockups based on descriptive prompts provided by the author. All generated designs were subsequently reviewed, refined, and adapted to meet the specific requirements of the NEST application.
\end{itemize}
All code, architectural designs, database schemas, and written analyses were reviewed, compiled, and verified by the author to ensure technical accuracy and academic integrity.